\scanpage{292}
\mspage{18}
% archive sheet number: 291

en notant que cette flèche passe au quotient en $\pi_0(X)\unsure{-|\longrightarrow}\mathrm{Hom}_{\unsure{W|w}}(e_W,X)$, vu que deux applications homotopes de $e_{\mathrm{Cat}}$ dans $X$ sont égales \unsure{dans} $\mathrm{Hot}_W$. Je dis que \deleted{\unsure{ce}}
\[
\beta\alpha=\mathrm{id},\qquad\alpha\beta=\mathrm{id}.
\]

La relation $\alpha\beta=\mathrm{id}$ revient à dire que la \deleted{l'application \illegible{1}}\inserted{flèche} $i_x:e_{\mathrm{Cat}}\unsure{-|\longrightarrow}X$ définie par $x\in\mathrm{Ob}\,\deleted{\unsure{X}}$, \deleted{donne un} donne $\pi_0(i_x)\deleted{\unsure{\text{élément}}}\inserted{{:}e_{\mathrm{Ens}}\to\pi_0(X)}$ \deleted{égale à} dont la valeur est la classe de $x$ dans $\pi_0(X)$, c'est tautologique. Pour prouver que $\beta\alpha=\mathrm{id}$, il suffit \inserted{\unsure{alors|donc}} de prouver que $\beta$ est surjective, i.e. le

\uline{\unsure{Korollaire|Corollaire}} 2.12. \inserted{Pour tt $X$ dans $\mathrm{Cat}$,} L'application
\[
\underbrace{\mathrm{Hom}_{\mathrm{Cat}}(e_{\mathrm{Cat}},X)}_{\simeq\mathrm{Ob}\,X}
\longrightarrow\mathrm{Hom}_W(e_W,X)
\]
est \uline{surjective} \deleted{\illegible{1}} (moyennant l'axiome de connexité Loc 5).

\deleted{Cela résulte du}

\uline{Dém. de 2.12}. Si $X=\varnothing$, l'assertion revient à dire que $\mathrm{Hom}_W(e_W,\varnothing)=\varnothing$, \deleted{\unsure{Or si}} \deleted{Mais il} mais plus généralement on a le

\uline{Cor. 2.1\deleted{\unsure{2}}3.} (Moyennant Loc 5) Soit $X$ non vide dans $\mathrm{Cat}$, alors $\mathrm{Hom}_W(X,\varnothing)=\varnothing$. \inserted{(Donc $\varnothing$ est un objet initial \uline{strict} de $\mathrm{Hot}_W$.)}

En effet, si on avait $X\to\varnothing$ dans $\mathrm{Hot}_W$,
